\section{Constraint algebra and degree-of-freedom count}\label{app:dirac}
The background-independent combined analysis uses a 42-dimensional phase space and a 24-constraint system. Six constraints are first class and eighteen are second class, giving
\begin{equation}
N_{\rm dof}=\frac{42-2N_{\rm FC}-N_{\rm SC}}{2}
=\frac{42-12-18}{2}=6.
\end{equation}
The second-class submatrix was verified to have rank 18 on deterministic generic algebraic patches, while the expected first-class null directions were identified separately. The nonseparable scalar operator adds neither a configuration variable nor higher time derivatives, and the VCDM auxiliary variables are removed by second-class/algebraic constraints rather than becoming an extra propagating scalar. The AeST part of this count is consistent with the independent fully nonlinear Hamiltonian analysis of Bataki, Skordis, and Zlosnik \cite{BatakiSkordisZlosnik2024}; the absence of an extra local scalar in VCDM is consistent with the Type-II minimally modified gravity analysis \cite{DeFeliceMaedaMukohyamaPookkillath2022}.

For local preferred-frame behavior we quote only the screened Einstein--aether limiting specialization used as a structural check. In that limit,
\begin{equation}
\alpha_1=-4K_B^{\rm loc},
\end{equation}
\begin{equation}
\alpha_2=\frac{K_B^{\rm loc}[K_B^{\rm loc}-c_2(1-2K_B^{\rm loc})]}{c_2(2-K_B^{\rm loc})}.
\end{equation}
The relation $c_2=K_B^{\rm loc}/(1-2K_B^{\rm loc})$ with $K_B^{\rm loc}=10^{-5}$ gives $\alpha_1=-4\times10^{-5}$ and $\alpha_2=0$ in this screened limit \cite{FosterJacobson2006}. This appendix does not claim that the limiting formula replaces a complete moving-source finite-field 1PN calculation of the full theory.
